\textbf{COLUNI 2023 -- Questão 02}\par\medskip “MINHA QUERIDA DONA, sei que você anda se queixando de mim, publicamente, até. Você não pode imaginar o sofrimento que isto me causa, mesmo porque você provavelmente acha que rádios são objetos inanimados, sem vida própria. Você está enganada. Ao menos no meu caso, você está enganada. Ao contrário do que você pensa, tenho sentimentos, tenho emoções. É em nome desses sentimentos e dessas emoções que lhe falo agora, tanto em AM como em FM. Na verdade, eu nem tinha tomado conhecimento de minha própria existência, até que fui instalado em seu carro”.

\par

SCLIAR, Moacyr. O rádio apaixonado. In: SCLIAR, Moacyr. \textbf{Histórias que os jornais não contam}. 3. ed. Porto Alegre: L\&M Editores, 2018. p.14. (Adaptado).

\par

No trecho acima, retirado da crônica “O rádio apaixonado”, há a predominância da figura de linguagem conhecida como:\begin{enumerate}
\item metáfora.
\item metonímia.
\item personificação.
\item comparação.\end{enumerate}
